\documentclass[../main]{subfiles}
\ifSubfilesClassLoaded{\setcounter{chapter}{9}}{}
\begin{document}

\chapter{UI Canvas dan Text}

\begin{subcpmk}
  \item Sub-CPMK 5.1: Membuat UI Canvas dan menampilkan skor
\end{subcpmk}

\section{Audio Source dan Audio Listener}

\begin{learningoutcomes}
\item Memahami Audio Clip, Audio Source, Audio Listener
\item Memainkan SFX dan BGM
\item Mengatur volume dan loop
\end{learningoutcomes}

\subsection{Komponen Audio}

Audio Clip: file suara (.mp3, .wav). Audio Source: pemancar suara di GameObject. Audio Listener: biasanya di Main Camera (satu saja per scene). PlayOnAwake, Loop, Volume.

\subsection{Implementasi SFX}

AudioSource.PlayOneShot(clip) untuk sound effect (tembakan, ledakan). Tidak menghentikan clip yang sedang diputar. Attach AudioSource ke Player atau buat prefab audio.

\section{BGM dan Audio Mixer}

\begin{learningoutcomes}
\item Memainkan Background Music
\item Menggunakan Audio Mixer untuk kontrol volume
\item Mengatur 3D vs 2D audio
\end{learningoutcomes}

\subsection{Background Music}

AudioSource dengan Loop = true. Play() di Start. DontDestroyOnLoad untuk BGM antar scene (opsional). Satu GameObject untuk BGM manager.

\subsection{Audio Mixer}

Window $\rightarrow$ Audio $\rightarrow$ Audio Mixer. Buat Groups: Master, Music, SFX. Expose parameter untuk kontrol dari script. Attenuation: 2D (UI, BGM) vs 3D (spatial, posisi matters).


\begin{aktivitas}
  \item Pelajari konsep Canvas, RectTransform, dan UI Text melalui materi di atas
  \item Diskusikan peran UI dalam game dan layout
  \item Analisis studi kasus Canvas dan penampilan skor
\end{aktivitas}

\begin{latihan}
  \item Apa fungsi RectTransform dan Anchor?
  \item Bagaimana UI Text dihubungkan ke script?
\end{latihan}

\begin{checklist}
  \item Saya dapat menjelaskan konsep Canvas dan UI
  \item Saya memahami RectTransform dan penampilan skor
\end{checklist}

\begin{rangkuman}
Bab ini membahas teori Canvas, UI Text, dan menampilkan skor.
\end{rangkuman}

\ifSubfilesClassLoaded{
  \renewcommand{\bibname}{Daftar Pustaka}
  \bibliographystyle{plain}
  \bibliography{../references}
}{}
\end{document}
